\section{Diagrama de caso de uso}

O diagrama de caso de uso é uma das partes fundamentais do projeto, pois é com ele que serão definidos os requisitos funcionais de forma clara do sistema. Na Figura~\ref{fig:caso-de-uso} são apresentados os principais atores e funcionalidades da aplicação.

\begin{figure}[H]
    \caption{Diagrama de caso de uso do sistema de reserva de salas}
    \label{fig:caso-de-uso}
    \centering
    \includegraphics[width=0.95\textwidth,height=0.75\textheight,keepaspectratio]{../../../out/src/plantuml/usecasediagram.png}

    {\footnotesize \textbf{Fonte:} O autor.
    \textbf{Acesso:} \href{\linkUsecasediagram}{Diagrama de Caso de Uso}.}
\end{figure}

\subsection{Atores}

\begin{description}
    \item[\textbf{sa\_professor}] Representa o usuário comum do sistema. Associa-se diretamente ao gerenciamento de solicitações, de ambientes e de agendamentos, além da geração de relatórios.

    \item[\textbf{sa\_coordenador}] Especialização de \textbf{sa\_professor}, herdando todas as suas permissões. Possui acesso adicional ao gerenciamento de usuários do sistema.
\end{description}

\subsection{Casos de uso}

\begin{description}
    \item[\textbf{suc\_gerenciar\_ambientes}] Reúne o cadastro, a edição, a remoção e a consulta dos elementos que compõem o espaço físico da instituição.

    \begin{description}
        \item[\textbf{suc\_gerenciar\_prédios}] Trata dos prédios, que delimitam onde cada sala está localizada.

        \item[\textbf{suc\_gerenciar\_salas}] Trata das salas propriamente ditas, que são as unidades efetivamente reserváveis.

        \item[\textbf{suc\_gerenciar\_tipos\_de\_salas}] Trata da classificação das salas, por exemplo laboratório, auditório ou sala de aula.

        \item[\textbf{suc\_gerenciar\_recursos}] Trata dos recursos associados às salas, como projetores e computadores.
    \end{description}

    \item[\textbf{suc\_gerenciar\_agendamentos}] Permite criar, atualizar, cancelar e visualizar agendamentos de salas.

    \item[\textbf{suc\_gerenciar\_solicitações}] Permite registrar e acompanhar as solicitações de reserva, com seus respectivos status.

    \item[\textbf{suc\_gerar\_relatórios}] Permite emitir relatórios com informações consolidadas sobre uso, agendamentos e histórico do sistema.

    \item[\textbf{suc\_exportar\_relatórios}] Estende \textbf{suc\_\allowbreak{}gerar\_\allowbreak{}relatórios} com a exportação do resultado em arquivo. É um comportamento opcional: o relatório pode ser apenas consultado, sem ser exportado.

    \item[\textbf{suc\_gerenciar\_usuarios}] Permite cadastrar, editar, remover e consultar usuários do sistema.
\end{description}
